\section{Индикаторы событий и ковариация}

\begin{definition}[Индикатор события]
\leavevmode

Индикатор события $A$:
$$ I_A = \begin{cases} 1, & \text{событие } A \text{ произошло} \\ 0, & \text{иначе} \end{cases} $$
\end{definition}

\begin{statement}[Свойства индикаторов]
\leavevmode

\begin{itemize}
    \item $E[I_A] = P(A)$
    \item $I_A \cdot I_B = I_{AB}$ (индикатор пересечения)
    \item Если $A$ и $B$ независимы, то $P(I_A = \alpha, I_B = \beta) = P(I_A = \alpha) \cdot P(I_B = \beta)$
\end{itemize}
\end{statement}

\begin{definition}[Ковариация]
\leavevmode

$$ \text{Cov}(\xi, \eta) = E(\xi \eta) - E(\xi) E(\eta) $$
\end{definition}

\begin{statement}
Независимые случайные величины некоррелированы ($\text{Cov} = 0$), но обратное неверно.
\end{statement}

\begin{example} Зависимые некоррелированные величины:
\begin{gather*}
X = \begin{pmatrix} -1 & 0 & 1 \\ 1/4 & 1/2 & 1/4 \end{pmatrix} \quad Y = X^2 \\
E(XY) = E(X^3) = (-1)^3 \cdot \frac{1}{4} + 0 \cdot \frac{1}{4} + 1^3 \cdot \frac{1}{4} = 0\\
E(X) = -1 \cdot \frac{1}{4} + 0 \cdot \frac{1}{4} + 1 \cdot \frac{1}{4} = 0\\
E(Y) = E(X^2) = (-1)^2 \cdot \frac{1}{4} + 0 \cdot \frac{1}{4} + 1^2 \cdot \frac{1}{4} = \frac{1}{2}\\
\text{Cov}(X, Y) = E(X Y) - E(X) E(Y) = 0
\end{gather*}

\end{example}


\begin{statement}[Дисперсия суммы]
$$ D(\xi + \eta) = D\xi + D\eta + 2\text{Cov}(\xi, \eta) $$
\end{statement}

\begin{proof}[Доказательство]
\leavevmode

\begin{gather*}
D(\xi + \eta) = E((\xi + \eta)^2) - (E(\xi + \eta))^2 = \\
= E(\xi^2 + 2 \xi \eta + \eta^2) - E\xi^2 - 2 (E\xi) (E\eta) - E\eta^2 = \\
= E(\xi^2) + 2 E(\xi \eta) + E(\eta^2) - E\xi^2 - 2 (E\xi) (E\eta) - E\eta^2 = \\
= \left(E(\xi^2) - E\xi^2\right) + \left(E(\eta^2) - E\eta^2\right) + 2 \left(E(\xi \eta) - (E\xi) (E\eta)\right)
\end{gather*}

По определению, первое слагаемое — это дисперсия $D\xi$, второе — дисперсия $D\eta$, а третье — удвоенная ковариация $2\text{Cov}(\xi, \eta)$:
 
$$D(\xi + \eta) = D\xi + D\eta + 2\text{Cov}(\xi, \eta)$$

\end{proof}

\vfill

% Для индикаторов:
% $$ \text{Cov}(I_A, I_B) = E(I_A I_B) - E(I_A) E(I_B) = P(AB) - P(A)P(B) $$

\begin{flushright}
    \textit{to be continued}
\end{flushright}